\section{Context Statistics}
\label{app:stats}

This appendix defines the context statistics used in Sections \ref{sec:predict} and \ref{sec:predictfull} as implemented in the accompanying code repository.

\subsection{Common measurement protocol}

Unless stated otherwise, every statistic is computed \emph{per channel} on the \emph{train split} of the raw series (the identical functions are evaluated on a temporary train-only dataset), restricted to the first $n_{\max} = 12000$ rows of that split, so no validation or test information enters any feature in this paper; for datasets with more than 50 channels we draw 50 channels without replacement using \texttt{numpy.default\_rng(0)} (datasets with fewer channels use all of them), and the reported dataset value is the mean across channels. ACF always denotes the Pearson correlation between $x_{1:n-\ell}$ and $x_{\ell+1:n}$ after subtracting the mean; a statistic returns 0 (or NaN for ACF) on degenerate constant inputs ($\mathrm{std} < 10^{-8}$).

\subsection{Definitions}

\paragraph{Drift index.} Split the series into $n_{\mathrm{seg}} = 20$ contiguous segments of (near-)equal length and measure the variance of the segment means relative to the total variance:
\begin{equation}
\mathrm{drift}(x) \;=\; \frac{\mathrm{Var}\big(\{\bar{x}_s\}_{s=1}^{n_{\mathrm{seg}}}\big)}{\mathrm{Var}(x) + 10^{-12}},
\qquad \bar{x}_s = \text{mean of segment } s .
\end{equation}
It is near 0 for stationary level processes and approaches 1 when level shifts dominate the total variance.

\paragraph{Long-lag autocorrelation.} ACF at a lag of one quarter of the (truncated) series length,
\begin{equation}
\mathrm{long\_acf}(x) \;=\; \mathrm{ACF}\big(x,\; \max(1, \lfloor 0.25\, n \rfloor)\big) .
\end{equation}

\paragraph{Detrended spectral periodicity.} Subtract the mean and a least-squares linear trend, take the power spectrum $P = |\mathrm{rFFT}(x_{\mathrm{dt}})|^2$, zero the DC bin $P_0$, and report the peak-to-total power ratio
\begin{equation}
\mathrm{periodicity\_dt}(x) \;=\; \frac{\max_k P_k}{\sum_k P_k + 10^{-12}} .
\end{equation}

\paragraph{Calendar-aligned autocorrelation (differenced).} Remove the mean and linear trend, take a first difference, $x' = \Delta\,\mathrm{detrend}(x - \bar{x})$, and average the ACF at known calendar lags $\mathcal{P}$:
\begin{equation}
\mathrm{calendar\_acf\_diff}(x) \;=\; \frac{1}{|\mathcal{P}|} \sum_{p \in \mathcal{P}} \mathrm{ACF}(x', p),
\end{equation}
with $\mathcal{P}$ per dataset in Table \ref{tab:calendar}. The companion $\mathrm{long\_acf\_diff}$ is the 25\%-lag ACF of the same detrended-and-differenced series $x'$.

\begin{table}[t]
\caption{Calendar period lags $\mathcal{P}$ (in timesteps) aligned to each dataset's sampling rate: the daily and weekly cycles, except exchange rate (business days: one week, one month).}
\label{tab:calendar}
\centering\small
\begin{tabular}{lll}
\toprule
Dataset & Sampling & $\mathcal{P}$ \\
\midrule
ETTh1, ETTh2, Electricity, Traffic & hourly & $[24, 168]$ \\
ETTm1, ETTm2 & 15-min & $[96, 672]$ \\
Weather & 10-min & $[144, 1008]$ \\
Exchange & daily & $[5, 30]$ \\
solar & 10-min & $[144, 1008]$ \\
PEMS04/08 & 5-min & $[288, 2016]$ \\
ERCOT & hourly & $[24, 168]$ \\
pedestrian & hourly & $[24, 168]$ \\
bikes & hourly & $[24, 168]$ \\
\bottomrule
\end{tabular}
\end{table}

\subsection{Measured values}

\begin{table}[t]
\caption{Dataset-level statistic values (channel means) on the eight main benchmarks (top) and the six extension datasets (bottom), measured on the train split under the protocol above. These are the features behind the predictability analysis of Section \ref{sec:predictfull}: drift correlates negatively and $\mathrm{calendar\_acf\_diff}$ positively with the long-context benefit.}
\label{tab:stats_values}
\centering\small
\begin{tabular}{lccccc}
\toprule
Dataset & periodicity\_dt & drift & long\_acf & calendar\_acf\_diff & long\_acf\_diff \\
\midrule
ETTh1 & 0.313 & 0.525 & 0.184 & 0.296 & 0.068 \\
ETTh2 & 0.270 & 0.598 & 0.023 & 0.251 & 0.023 \\
ETTm1 & 0.181 & 0.523 & 0.122 & 0.106 & $-0.005$ \\
ETTm2 & 0.171 & 0.542 & 0.153 & 0.130 & $-0.013$ \\
Exchange & 0.478 & 0.897 & $-0.259$ & 0.001 & 0.011 \\
Weather & 0.146 & 0.275 & 0.048 & 0.061 & 0.001 \\
Electricity & 0.477 & 0.181 & 0.540 & 0.683 & 0.290 \\
Traffic & 0.475 & 0.027 & 0.627 & 0.510 & 0.253 \\
\midrule
solar & 0.364 & 0.033 & 0.176 & 0.234 & $-$0.026 \\
PEMS04 & 0.163 & 0.090 & $-$0.101 & 0.096 & 0.001 \\
PEMS08 & 0.225 & 0.071 & $-$0.261 & 0.163 & $-$0.004 \\
ERCOT & 0.270 & 0.340 & 0.106 & 0.900 & 0.410 \\
pedestrian & 0.506 & 0.010 & 0.697 & 0.649 & 0.151 \\
bikes & 0.202 & 0.016 & 0.371 & 0.258 & 0.041 \\
\bottomrule
\end{tabular}
\end{table}

\subsection{Lessons from the first version (why detrend/difference is mandatory)}

An earlier pilot version of our statistics measured periodicity by the \emph{raw} spectral peak ratio and autocorrelation on the \emph{raw} series, and both were contaminated by drift. A strong trend concentrates low-frequency power in the spectrum, where it masquerades as a periodic peak: exchange rate scored a spurious periodicity of $0.558$ in v1 despite having no meaningful periodic structure (its detrended score is $0.478$, and its calendar-lag ACF after differencing is $\approx 0$). Symmetrically, raw long-lag ACF is inflated by level drift --- a slowly wandering level keeps the series correlated with itself at any lag --- so calendar and long-lag autocorrelations are only meaningful after detrending \emph{and} first differencing. All statistics reported in the paper use the v2 definitions above.

\paragraph{Extension-suite statistics (v3).}
On the enlarged suite (fourteen datasets, twenty-eight samples), the drift index stays the dominant statistic: pooled over the twenty-eight samples it correlates at $-0.654$ with the log-span ($p = 2\times10^{-4}$) and $-0.768$ with the long-context benefit ($p < 10^{-5}$), and at the dataset level ($n{=}14$, the statistically independent units) at $-0.80$ with the log-span ($p = 5\times10^{-4}$) and $-0.85$ with the benefit ($p = 1\times10^{-4}$). The calendar-aligned differenced ACF is directionally consistent but weaker than the pooled view suggested: $+0.366$ ($p=0.056$) over twenty-eight samples, but the two samples per dataset share identical feature values, so the effective units are fourteen; at $n{=}14$ the correlation is $+0.425$ with $p=0.130$, and it does not survive Holm correction across the ten statistic$\times$target tests that Appendix \ref{app:bocpd} applies by the same standard. We therefore report calendar as a supporting descriptor, not an independent axis of significance. Long-lag ACF and detrended spectral periodicity drop below significance because their earlier signal was carried by the electricity and traffic extremes. Full per-dataset values, source notes, and the updated correlation tables are released in the accompanying repository.

\paragraph{Drift-index conventions.}
An earlier version of this analysis computed the statistics on the first 12000 rows of the whole file, which for the shorter datasets includes validation and test spans. The two conventions differ only where the first 12000 rows cross the train boundary (ETTh1 0.450 full-series vs.\ 0.525 train-split; ETTh2 0.651 vs.\ 0.598; exchange rate 0.894 vs.\ 0.897); every dataset whose train split exceeds 12000 rows is identical under both. Switching to the leakage-free train-split convention leaves every correlation direction and the dataset-level significance pattern intact, and the numbers reported in this paper use it throughout; the residual difference matters only near the veto threshold of the Chronos budget rule, which we disclose in Appendix \ref{app:select}.
